% --- Archon physics formalization source begin ---
% archon:physics
% archon:covers PhyXMiniProblems/problem_phyx_mini_0141.lean
% archon:source-report reports/phyx_mini/problem_phyx_mini_0141.source.json

\chapter{Physics problem phyx\_mini\_0141}
\label{ch:PhyXMiniProblems_problem_phyx_mini_0141}

\paragraph{Problem source.}
\#\# Physical scenario

A woman \textbackslash{}(1.60\textbackslash{},\textbackslash{}mathrm\{m\}\textbackslash{}) tall stands in front of a vertical plane mirror.She is to be able to see her whole body. Assume that her eyes are \textbackslash{}(10\textbackslash{},\textbackslash{}mathrm\{cm\}\textbackslash{}) below the top of her head.

\#\# Question

How high must the mirror's lower edge be above the floor?

\#\# Answer choices

- A: A: 0.85m

- B: B: 0.80m

- C: C: 0.75m

- D: D: 0.70m

\#\# Figure caption (auxiliary; use the image as primary evidence)

The image illustrates the concept of reflection in a plane mirror. It features a side view of a person standing in front of a vertical mirror. 

- **Person and Image**: 
  - The person, holding a briefcase in their right hand, is positioned on the left side of the mirror.
  - The image, shown on the mirror’s right side, mirrors the person's posture and appearance.

- **Mirror**:
  - The mirror is depicted as a vertical line labeled "Mirror". 

- **Height Measurements**:
  - The person is shown with a height of 1.50 m. 
  - An additional measurement indicates a 0.10 m distance from the top of the head to point G.

- **Lines and Points**:
  - Several pink lines represent light rays:
    - AE and AG: Rays directed from the person to the mirror.
    - EF and GH: Reflected rays that meet at the mirror, showing path alignments.
    - BC: Represents the reflected image of the person.

- **Marked Points**:
  - A, B, C, D, E, F, G, and H are marked points or intersections related to the person's position, the mirror, and the reflected rays.

- **Annotations**:
  - "Image" label is placed near the right to denote the mirrored image of the person.

The diagram effectively demonstrates the principles of reflection and image formation in a flat mirror.

\#\# Dataset metadata

Geometrical Optics; Spatial Relation Reasoning, Physical Model Grounding Reasoning

\paragraph{Recorded answer/context.}
C: C: 0.75m

\paragraph{Figure/image path.}
/root/proposal\_for\_physic/science-mango/phyx\_mini\_run/phyx\_data/test\_image/141.png

\paragraph{Formalization target.}
create a compiling Lean file with sorry bodies at `PhyXMiniProblems/problem\_phyx\_mini\_0141.lean`. The Lean declarations must preserve the physical quantities, dimensions or dimensional roles, figure labels, governing-law hypotheses, and final relation expressed by this problem.
Use Mathlib/Physlib names found through LeanExplore where available. If a physics API is missing, introduce faithful local abstractions rather than scalar placeholder aliases.

\begin{theorem}[Physics formalization target]
\label{thm:physics:phyx_mini_0141:target}
\lean{PhyXMiniProblems.ProblemPhyXMini0141.lowerMirrorEdgeHeight_eq_threeQuartersMeter}
\uses{def:physics:phyx-mini-0141:phyxminiproblems-problemphyxmini0141-fullbodyplanemirrordiagram, def:physics:phyx-mini-0141:phyxminiproblems-problemphyxmini0141-verticalseparationinmeters, def:physics:phyx-mini-0141:phyxminiproblems-problemphyxmini0141-matchesproblemreadouts, def:physics:phyx-mini-0141:phyxminiproblems-problemphyxmini0141-matchesprimaryfigure, def:physics:phyx-mini-0141:phyxminiproblems-problemphyxmini0141-obeysidealplanemirroroptics}
The assigned autoformalize agent should translate this physics problem into Lean declarations in the covered file, with theorem and lemma proof bodies written as `by sorry`.
\end{theorem}
\begin{proof}
This is an autoformalization task, not a proof task. Produce statements that can later be proved by the physics prover without weakening the physical model.
\end{proof}

% --- Archon declaration topology begin ---
\section{Declaration topology}
\begin{definition}[Length Quantity]
  \label{def:physics:phyx-mini-0141:phyxminiproblems-problemphyxmini0141-lengthquantity}
  \lean{PhyXMiniProblems.ProblemPhyXMini0141.LengthQuantity}
  A signed physical length, independent of the unit used for its readout
\end{definition}
\begin{proof}
  This is the declared model definition of Length Quantity.
\end{proof}
\begin{definition}[metre Unit Choices]
  \label{def:physics:phyx-mini-0141:phyxminiproblems-problemphyxmini0141-metreunitchoices}
  \lean{PhyXMiniProblems.ProblemPhyXMini0141.metreUnitChoices}
  SI units, with the length component made explicitly equal to metres
\end{definition}
\begin{proof}
  This is the declared model definition of metre Unit Choices.
\end{proof}
\begin{definition}[length In Meters]
  \label{def:physics:phyx-mini-0141:phyxminiproblems-problemphyxmini0141-lengthinmeters}
  \lean{PhyXMiniProblems.ProblemPhyXMini0141.lengthInMeters}
  \uses{def:physics:phyx-mini-0141:phyxminiproblems-problemphyxmini0141-lengthquantity, def:physics:phyx-mini-0141:phyxminiproblems-problemphyxmini0141-metreunitchoices}
  The scalar readout of a physical length in metres
\end{definition}
\begin{proof}
  This is the declared model definition of length In Meters.
\end{proof}
\begin{definition}[Diagram Point]
  \label{def:physics:phyx-mini-0141:phyxminiproblems-problemphyxmini0141-diagrampoint}
  \lean{PhyXMiniProblems.ProblemPhyXMini0141.DiagramPoint}
  \uses{def:physics:phyx-mini-0141:phyxminiproblems-problemphyxmini0141-lengthquantity}
  A physical point in the vertical plane of the person and mirror
\end{definition}
\begin{proof}
  This is the declared model definition of Diagram Point.
\end{proof}
\begin{definition}[Figure Point]
  \label{def:physics:phyx-mini-0141:phyxminiproblems-problemphyxmini0141-figurepoint}
  \lean{PhyXMiniProblems.ProblemPhyXMini0141.FigurePoint}
  The eight point labels visible in the source figure
\end{definition}
\begin{proof}
  This is the declared model definition of Figure Point.
\end{proof}
\begin{definition}[Figure Ray Part]
  \label{def:physics:phyx-mini-0141:phyxminiproblems-problemphyxmini0141-figureraypart}
  \lean{PhyXMiniProblems.ProblemPhyXMini0141.FigureRayPart}
  Solid ray pieces and dashed virtual extensions drawn in the primary image
\end{definition}
\begin{proof}
  This is the declared model definition of Figure Ray Part.
\end{proof}
\begin{definition}[start Point]
  \label{def:physics:phyx-mini-0141:phyxminiproblems-problemphyxmini0141-figureraypart-startpoint}
  \lean{PhyXMiniProblems.ProblemPhyXMini0141.FigureRayPart.startPoint}
  \uses{def:physics:phyx-mini-0141:phyxminiproblems-problemphyxmini0141-figurepoint, def:physics:phyx-mini-0141:phyxminiproblems-problemphyxmini0141-figureraypart}
  Initial endpoint of each labeled ray part in the figure
\end{definition}
\begin{proof}
  This is the declared model definition of start Point.
\end{proof}
\begin{definition}[end Point]
  \label{def:physics:phyx-mini-0141:phyxminiproblems-problemphyxmini0141-figureraypart-endpoint}
  \lean{PhyXMiniProblems.ProblemPhyXMini0141.FigureRayPart.endPoint}
  \uses{def:physics:phyx-mini-0141:phyxminiproblems-problemphyxmini0141-figurepoint, def:physics:phyx-mini-0141:phyxminiproblems-problemphyxmini0141-figureraypart}
  Final endpoint of each labeled ray part in the figure
\end{definition}
\begin{proof}
  This is the declared model definition of end Point.
\end{proof}
\begin{definition}[Full Body Plane Mirror Diagram]
  \label{def:physics:phyx-mini-0141:phyxminiproblems-problemphyxmini0141-fullbodyplanemirrordiagram}
  \lean{PhyXMiniProblems.ProblemPhyXMini0141.FullBodyPlaneMirrorDiagram}
  \uses{def:physics:phyx-mini-0141:phyxminiproblems-problemphyxmini0141-diagrampoint, def:physics:phyx-mini-0141:phyxminiproblems-problemphyxmini0141-figurepoint, def:physics:phyx-mini-0141:phyxminiproblems-problemphyxmini0141-figureraypart}
  The physical point positions and the ray parts shown in the diagram
\end{definition}
\begin{proof}
  This is the declared model definition of Full Body Plane Mirror Diagram.
\end{proof}
\begin{definition}[horizontal Position In Meters]
  \label{def:physics:phyx-mini-0141:phyxminiproblems-problemphyxmini0141-horizontalpositioninmeters}
  \lean{PhyXMiniProblems.ProblemPhyXMini0141.horizontalPositionInMeters}
  \uses{def:physics:phyx-mini-0141:phyxminiproblems-problemphyxmini0141-lengthinmeters, def:physics:phyx-mini-0141:phyxminiproblems-problemphyxmini0141-figurepoint, def:physics:phyx-mini-0141:phyxminiproblems-problemphyxmini0141-fullbodyplanemirrordiagram}
  Horizontal metre coordinate of a labeled point
\end{definition}
\begin{proof}
  This is the declared model definition of horizontal Position In Meters.
\end{proof}
\begin{definition}[vertical Position In Meters]
  \label{def:physics:phyx-mini-0141:phyxminiproblems-problemphyxmini0141-verticalpositioninmeters}
  \lean{PhyXMiniProblems.ProblemPhyXMini0141.verticalPositionInMeters}
  \uses{def:physics:phyx-mini-0141:phyxminiproblems-problemphyxmini0141-lengthinmeters, def:physics:phyx-mini-0141:phyxminiproblems-problemphyxmini0141-figurepoint, def:physics:phyx-mini-0141:phyxminiproblems-problemphyxmini0141-fullbodyplanemirrordiagram}
  Vertical metre coordinate of a labeled point
\end{definition}
\begin{proof}
  This is the declared model definition of vertical Position In Meters.
\end{proof}
\begin{definition}[point In Meters]
  \label{def:physics:phyx-mini-0141:phyxminiproblems-problemphyxmini0141-pointinmeters}
  \lean{PhyXMiniProblems.ProblemPhyXMini0141.pointInMeters}
  \uses{def:physics:phyx-mini-0141:phyxminiproblems-problemphyxmini0141-figurepoint, def:physics:phyx-mini-0141:phyxminiproblems-problemphyxmini0141-fullbodyplanemirrordiagram, def:physics:phyx-mini-0141:phyxminiproblems-problemphyxmini0141-horizontalpositioninmeters, def:physics:phyx-mini-0141:phyxminiproblems-problemphyxmini0141-verticalpositioninmeters}
  Two-dimensional metre readout used only for affine diagram geometry
\end{definition}
\begin{proof}
  This is the declared model definition of point In Meters.
\end{proof}
\begin{definition}[vertical Separation In Meters]
  \label{def:physics:phyx-mini-0141:phyxminiproblems-problemphyxmini0141-verticalseparationinmeters}
  \lean{PhyXMiniProblems.ProblemPhyXMini0141.verticalSeparationInMeters}
  \uses{def:physics:phyx-mini-0141:phyxminiproblems-problemphyxmini0141-figurepoint, def:physics:phyx-mini-0141:phyxminiproblems-problemphyxmini0141-fullbodyplanemirrordiagram, def:physics:phyx-mini-0141:phyxminiproblems-problemphyxmini0141-verticalpositioninmeters}
  Signed vertical separation, in metres, from lower to upper
\end{definition}
\begin{proof}
  This is the declared model definition of vertical Separation In Meters.
\end{proof}
\begin{definition}[Matches Problem Readouts]
  \label{def:physics:phyx-mini-0141:phyxminiproblems-problemphyxmini0141-matchesproblemreadouts}
  \lean{PhyXMiniProblems.ProblemPhyXMini0141.MatchesProblemReadouts}
  \uses{def:physics:phyx-mini-0141:phyxminiproblems-problemphyxmini0141-fullbodyplanemirrordiagram, def:physics:phyx-mini-0141:phyxminiproblems-problemphyxmini0141-verticalseparationinmeters}
  The numerical measurements stated in the problem. In particular, this predicate contains no information about the height of point B
\end{definition}
\begin{proof}
  This is the declared model definition of Matches Problem Readouts.
\end{proof}
\begin{definition}[Matches Primary Figure]
  \label{def:physics:phyx-mini-0141:phyxminiproblems-problemphyxmini0141-matchesprimaryfigure}
  \lean{PhyXMiniProblems.ProblemPhyXMini0141.MatchesPrimaryFigure}
  \uses{def:physics:phyx-mini-0141:phyxminiproblems-problemphyxmini0141-fullbodyplanemirrordiagram, def:physics:phyx-mini-0141:phyxminiproblems-problemphyxmini0141-horizontalpositioninmeters, def:physics:phyx-mini-0141:phyxminiproblems-problemphyxmini0141-verticalpositioninmeters, def:physics:phyx-mini-0141:phyxminiproblems-problemphyxmini0141-verticalseparationinmeters}
  Incidence, alignment, ordering, and calibrated readouts visible in the primary image. The 1.50 m arrow is the floor-to-eye segment; together with the 0.10 m arrow above it, it agrees with the stated total height 1.60 m. No numerical position for the lower mirror edge is recorded here
\end{definition}
\begin{proof}
  This is the declared model definition of Matches Primary Figure.
\end{proof}
\begin{definition}[Is Virtual Image Across Vertical Mirror]
  \label{def:physics:phyx-mini-0141:phyxminiproblems-problemphyxmini0141-isvirtualimageacrossverticalmirror}
  \lean{PhyXMiniProblems.ProblemPhyXMini0141.IsVirtualImageAcrossVerticalMirror}
  \uses{def:physics:phyx-mini-0141:phyxminiproblems-problemphyxmini0141-figurepoint, def:physics:phyx-mini-0141:phyxminiproblems-problemphyxmini0141-fullbodyplanemirrordiagram, def:physics:phyx-mini-0141:phyxminiproblems-problemphyxmini0141-horizontalpositioninmeters, def:physics:phyx-mini-0141:phyxminiproblems-problemphyxmini0141-verticalpositioninmeters}
  The virtual-image construction for a vertical plane mirror: image and object have equal height and equal, opposite horizontal displacement from the mirror. This is a governing plane-mirror law and does not prescribe the height of the mirror point
\end{definition}
\begin{proof}
  This is the declared model definition of Is Virtual Image Across Vertical Mirror.
\end{proof}
\begin{definition}[Obeys Ideal Plane Mirror Optics]
  \label{def:physics:phyx-mini-0141:phyxminiproblems-problemphyxmini0141-obeysidealplanemirroroptics}
  \lean{PhyXMiniProblems.ProblemPhyXMini0141.ObeysIdealPlaneMirrorOptics}
  \uses{def:physics:phyx-mini-0141:phyxminiproblems-problemphyxmini0141-fullbodyplanemirrordiagram, def:physics:phyx-mini-0141:phyxminiproblems-problemphyxmini0141-pointinmeters, def:physics:phyx-mini-0141:phyxminiproblems-problemphyxmini0141-isvirtualimageacrossverticalmirror}
  Ideal plane-mirror optics for the two boundary rays. The reflected ray to the eye is collinear with its dashed continuation toward the corresponding virtual image. Mathlib's Collinear supplies the affine-geometric notion; the equal-distance clauses supply the physical plane-mirror image law
\end{definition}
\begin{proof}
  This is the declared model definition of Obeys Ideal Plane Mirror Optics.
\end{proof}
\begin{lemma}[eye Height eq three Halves]
  \label{lem:physics:phyx-mini-0141:phyxminiproblems-problemphyxmini0141-eyeheight-eq-threehalves}
  \lean{PhyXMiniProblems.ProblemPhyXMini0141.eyeHeight_eq_threeHalves}
  \uses{def:physics:phyx-mini-0141:phyxminiproblems-problemphyxmini0141-fullbodyplanemirrordiagram, def:physics:phyx-mini-0141:phyxminiproblems-problemphyxmini0141-verticalseparationinmeters, def:physics:phyx-mini-0141:phyxminiproblems-problemphyxmini0141-matchesproblemreadouts}
  The problem measurements imply that the eye is 1.50 m above the foot
\end{lemma}
\begin{proof}
  The conclusion follows from the governing relations and figure readouts named in the statement of eye Height eq three Halves.
\end{proof}
\begin{lemma}[lower Edge bisects eye To Floor height]
  \label{lem:physics:phyx-mini-0141:phyxminiproblems-problemphyxmini0141-loweredge-bisects-eyetofloor-height}
  \lean{PhyXMiniProblems.ProblemPhyXMini0141.lowerEdge_bisects_eyeToFloor_height}
  \uses{def:physics:phyx-mini-0141:phyxminiproblems-problemphyxmini0141-fullbodyplanemirrordiagram, def:physics:phyx-mini-0141:phyxminiproblems-problemphyxmini0141-verticalseparationinmeters, def:physics:phyx-mini-0141:phyxminiproblems-problemphyxmini0141-matchesprimaryfigure, def:physics:phyx-mini-0141:phyxminiproblems-problemphyxmini0141-obeysidealplanemirroroptics}
  At the limiting placement for seeing the feet, the lower edge B lies halfway in height between the eye E and the floor. This is a consequence of the vertical-mirror image law and the straight virtual sightline, not a setup assumption
\end{lemma}
\begin{proof}
  The conclusion follows from the governing relations and figure readouts named in the statement of lower Edge bisects eye To Floor height.
\end{proof}
\begin{definition}[Answer Choice]
  \label{def:physics:phyx-mini-0141:phyxminiproblems-problemphyxmini0141-answerchoice}
  \lean{PhyXMiniProblems.ProblemPhyXMini0141.AnswerChoice}
  The four displayed answer labels
\end{definition}
\begin{proof}
  This is the declared model definition of Answer Choice.
\end{proof}
\begin{definition}[lower Edge Height In Meters]
  \label{def:physics:phyx-mini-0141:phyxminiproblems-problemphyxmini0141-answerchoice-loweredgeheightinmeters}
  \lean{PhyXMiniProblems.ProblemPhyXMini0141.AnswerChoice.lowerEdgeHeightInMeters}
  \uses{def:physics:phyx-mini-0141:phyxminiproblems-problemphyxmini0141-answerchoice}
  Lower-edge height printed beside each answer choice, in metres
\end{definition}
\begin{proof}
  This is the declared model definition of lower Edge Height In Meters.
\end{proof}
\begin{definition}[recorded Answer Choice]
  \label{def:physics:phyx-mini-0141:phyxminiproblems-problemphyxmini0141-recordedanswerchoice}
  \lean{PhyXMiniProblems.ProblemPhyXMini0141.recordedAnswerChoice}
  \uses{def:physics:phyx-mini-0141:phyxminiproblems-problemphyxmini0141-answerchoice}
  Dataset metadata: the recorded answer is choice C
\end{definition}
\begin{proof}
  This is the declared model definition of recorded Answer Choice.
\end{proof}
% --- Archon declaration topology end ---
% --- Archon physics formalization source end ---
